\begin{helper}
For every \(2k\)-tuple
\[
  ab=(a_1,b_1,\ldots,a_k,b_k)\in(\mathbb F_2^p)^{2k}
\]
and every index \(i<k\), let \(r_i\) be the rank quantity from
\(\noderef{F2BadTupleRank}\), and let
\[
  U_i=\operatorname{span}\{a_j:j<i\}.
\]
If \(r_{i+1}=r_i\), then \(a_i\in U_i\), and
\[
  |U_i|\cdot
  \left|\{\,y\in\mathbb F_2^p:\langle a_j,y\rangle=1
      \text{ for every }j<i+1\,\}\right|
  \le
  2^{r_i}\,2^{p-r_i}.
\]
\end{helper}
\begin{proof}
The span defining \(r_{i+1}\) is obtained from the span defining \(r_i\) by
adjoining the single vector \(a_i\). If \(a_i\notin U_i\), adjoining \(a_i\)
adds one independent direction, so the dimension would increase by one:
\(r_{i+1}=r_i+1\). This contradicts the hypothesis \(r_{i+1}=r_i\). Hence
\(a_i\in U_i\).

By \(\noderef{F2BadTuplePrefixSpanCard}\), the span \(U_i\) has cardinality
\(2^{r_i}\). By \(\noderef{F2BadTuplePrefixFiberBound}\), applied to the prefix
of length \(i+1\), the number of vectors \(y\) satisfying
\(\langle a_j,y\rangle=1\) for every \(j<i+1\) is at most
\(2^{p-r_{i+1}}\). Since \(r_{i+1}=r_i\), this upper bound is
\(2^{p-r_i}\). Multiplying the two bounds gives the stated product estimate.
\end{proof}
